\section{Weight finiteness on the same deformation space}

The propagation theorem uses a radical containment on a matrix
deformation space. Arithmetic applications usually supply module
finiteness over a weight algebra on a different ring. We explain the
passage between them, and how a good prime on the matrix space
returns an actual determinant point.

Fix an original setup $T$ of rank $n$ (two or three), with coefficient
DVR $O$, residue field $k$ of order $q$, and universal ordinary
determinant ring $B^{\rm det}$. The matrix space $S=T.S$ has
matrix algebra of coordinates $A=S.A$ and its coefficient subring
$B=S.B$ of degree zero. Write $\Phi:B^{\rm det}\to B$ for an
original comparison map $T.\mathrm{IsPhi}$. Thus it represents the
determinant and the ordered characters of the universal matrices
after the specified scalar twist. Write $\lambda$ for $S.\mathrm{lamId}$
and $\mathfrak m_B$ for the contraction of the maximal ideal of
the complete universal matrix ring to $B$.

% BEGIN W67F1
\paragraph{W67F1. Specializing inertia in a quotient with nilpotents.}
Let $R$ be a complete Noetherian local $O$-algebra in the original
coefficient category, with identified residue field $k$. Let $I$
be any proper ideal containing the image of $\mathfrak m_O$. The
induced map $k\to R/I$ will be denoted by $j$. If a unit $y\in R$
satisfies $y^{q-1}=1\bmod I$, then
\[
                         y\bmod I=j(\bar y).
\]
Indeed, lift $\bar y$ to a unit $c\in O$. Then $c^{q-1}=1\bmod I$.
Factoring the difference of powers in $R$ gives $(y-c)$ times a
factor whose residue is $(q-1)\bar y^{q-2}\ne0$. This factor is
a unit, so $y-c\in I$. The proof works for $q=2$ and uses no
reducedness assumption on the quotient.

Let $\Lambda$ represent the original inertial weight characters
and let $f:\Lambda\to R$ be the map attached to an ordinary
determinant pair. Suppose every value $y$ of each of its ordered
characters on the prescribed inertia groups satisfies the power
condition above. The two $O$-maps $\Lambda\to R/I$ obtained from
$f$ and from $\Lambda\to k\xrightarrow{j}R/I$ then agree on all
universal inertia characters. The quotient is again in the original
coefficient category. Indeed it is a proper quotient of a complete
Noetherian local ring with the same residue field. The pushed
characters have the same reduction and prescribed central values,
and are continuous since a local homomorphism takes the $a$th
maximal-ideal power into the corresponding $a$th power.
The representing property for the inertial character functor
therefore makes the two maps equal. In particular
$f(\mathfrak m_\Lambda)\subset I$.
% END W67F1

% BEGIN W67F2
\paragraph{W67F2. Transport from module finiteness to the original weight ideal.}
Assume $\Phi$ is surjective and take any ideal $J\subset B^{\rm det}$
and original weight object $\Lambda$ for which $B^{\rm det}/J$
is finite over $\Lambda$ via the universal pair. Put $M=JB$.
We prove
\[
                 \mathfrak m_B\subset\sqrt{M+\lambda}.
\]
For a prime $Q\supset M+\lambda$ put $P=\Phi^{-1}(Q)$. It
contains $J$ and $\mathfrak m_OB^{\rm det}$. For each local inertia
element $g$, the image of the ordered determinant character value
under $\Phi$ is the product of the matrix-space character at the
linked local element and a scalar unit from the twist. The
$(q-1)$st power of the linked element is still in inertia and its
residual character $\Psi$ is trivial, since $k^\times$ has order
$q-1$. By the defining generators of $\lambda$, the first factor
to that power is one modulo $Q$. The scalar unit has the same
property because $\mathfrak m_O B\subset Q$. Thus all these
character values to the power $q-1$ are one modulo $P$.

W67F1 applies to $B^{\rm det}/P$ through the universal weight map
and gives $\mathfrak m_\Lambda\subset f^{-1}(P)$. The latter is
proper, so equality holds. The domain $B^{\rm det}/P$ is finite
over the residue field of $\Lambda$, being a quotient of
$B^{\rm det}/J$ with that action. It is therefore a field and
$P$ is maximal. Since $\Phi$ is onto, $B/Q\simeq B^{\rm det}/P$;
in particular $Q$ is maximal. The ring $B$ is local as a quotient
of $B^{\rm det}$, and $\mathfrak m_B$ is proper, so
$\mathfrak m_B\subset Q$. Intersect over all such primes. This is
the full source implication $T.\mathrm{FinLamOf}$ for the actual
map and ideal. It retains the required surjectivity premise.
% END W67F2

% BEGIN W67F3
\paragraph{W67F3. A good component supplies the determinant point.}
Here $\Phi$ may be any comparison map. Assume
$\mathrm{AllGood}(S,JB)$. Fix an original coefficient-domain
test $O\to A_0$, a lift $r$ with its prescribed diagonal reduction,
and an embedding $i:A_0\hookrightarrow K$ into a field over which
$r$ is irreducible. Use the subring
$C=\operatorname{im}(O)+\mathfrak m_{A_0}$ inside $A_0$, the
preimage of $k$ in its residue field. It is local with maximal ideal
$\mathfrak m_{A_0}$: the inverse of a unit with residue in
$k^\times$ again has residue there. Lifts of a residue-field basis
containing $1$ generate $A_0$ as a $C$-module, since the error
after reduction belongs to $\mathfrak m_{A_0}\subset C$.
As $A_0$ is Noetherian, the finite-module descent theorem gives
Noetherianity of $C$. The powers of the common maximal ideal agree
in the two rings (an ambient multiplier can be absorbed in one
factor). The preimage of $k$ is closed, so $C$ is complete.
Thus $C$ is an allowed coefficient ring with residue field $k$,
in either characteristic. Diagonal residual reduction puts every
matrix entry and every local character value of $r$ in $C$,
and inverses of character values also lie there. Continuity and
the defining normalized matrix identities restrict to $C$.
Universality of the complete matrix ring, restriction to $S.A$,
and the inclusion into $A_0$ give an $O$-map $g:S.A\to A_0$
realizing $r$ and its ordered characters in their original basis.
Put $p=\ker(i g)$. This is a prime. The type-I test
on $S$ applies to $p$ because the field realization is the
specified irreducible $r\otimes_i K$. Hence $JB\subset p\cap B$.
The map
\[
                B^{\rm det}\xrightarrow{\Phi}B
                          \hookrightarrow S.A\xrightarrow{g}A_0
\]
kills $J$ by injectivity of $i$.

It remains to check what this map represents, since containment
of $J$ alone only gives a factorization of rings. The inclusion
of $S.A$ in the complete matrix ring sends the universal matrix
and its ordered characters to the defining ones of $S$.
After that inclusion the determinant law and the characters agree
by $T.\mathrm{IsPhi}$. They already agree in $S.A$: equality of
the determinant polynomial laws can be checked on their coefficients
in polynomial rings, where extending an injective map stays
injective. The character equalities can be tested value by value.
Now apply $g$. Forming the determinant of the matrix representation,
multiplying it by the specified scalar twist before taking its
determinant, and pushing forward its ordered characters all commute
with coefficient maps. We obtain $T.\mathrm{GoodPt}(J)$ with
precisely $T.K.\mathrm{tw}(r)$ and $T.K.\mathrm{tw}\psi(r)$.
This proves the native $T.\mathrm{GoodPtOf}$.
% END W67F3

% BEGIN W67F4
\paragraph{W67F4. A finite-point test on the marked upper frame.}
Consider two original setups on the same frame, in ranks two and
three, with their respective surjective comparison maps and ideals
$J_2,J_3$. Suppose the two determinant quotients are finite over
their respective weight objects. Write $M_i=J_iS_i.B$. Suppose
$\mathrm{P1}'$ and $\mathrm{SeedIn}$ hold for $(S_2,M_2)$, and
$\mathrm{P1}',\mathrm{P23}',\mathrm{SeedIn}$ for $(S_3,M_3)$.
These refer to the indicated ideals and include the actual seed primes.

W67F2 supplies the two finite-weight conditions. Applying ACC26 in
rank two yields $\mathrm{AllGood}(S_2,M_2)$. Apply U67I5 in rank
three with this very rank-two space and ideal; it gives
\[
  \mathrm{AllGood}(S_3,M_3)\qquad\text{and}\qquad\mathrm{OffFin}(S_3).
\]
W67F3 returns both arithmetic good-point conclusions with $J_i$ in
their original determinant rings. ACP11 gives $\mathrm{FinPt}$
in rank three, for every test ring and every weight-zero point.
If the setups are on the extended coefficient frame of the marked
argument, RTC24 and the already proved return RTC25 descend this
conclusion to the original frame and yield the finite end with the
same mark. No new mark is chosen. Thus the finite-point side now has
a sufficient condition stated in terms of the two arithmetic setups
and their prime and seed conditions. Realizing those inputs on the
specific frame, and the separate upper automorphy and descent used
by M65A4, are still required for $\mathrm{THM\_S6\_S}$.
% END W67F4
